\chapter{Conclusion}
\label{chap:conclusion}

In this thesis, we have presented the importance of indirect dark matter searches, design, development, and testing of two balloon-borne experiments, \ac{GAPS} and \ac{GRAMS}, aimed at advancing our understanding of cosmic rays and dark matter. The \ac{GAPS} experiment successfully launched its science payload on a long-duration balloon flight, demonstrating the viability of its novel detection technique. The data collected from \ac{GAPS} are undergoing analysis, and we anticipate that they will provide valuable insights into the nature of cosmic rays and their interactions. GAPS expect to deliver the first precision measurement of the cosmic antiproton spectrum below 0.25 GeV, see Fig.~\ref{fig:GAPS_antiproton} and the first-ever measurement of the cosmic antideuteron spectrum in $<0.25 GerV/N$ energy range with a sensitivity of $2\times10^{-6}$ [m$^2$ s sr GeV/n]$^{-1}$, see Fig.~\ref{fig:GAPS_antideuteron}.


\begin{figure}
    \centering
    \includegraphics[width=1\textwidth]{fig/GAPS_antiproton_sensitivity.png}
    \noindent\raggedright The projected \ac{GAPS} precision cosmic antiproton spectrum (red) at the top of the atmosphere is shown with the statistics expected from three 35-day flights \cite{ROGERS2023102791}. Data from BESS (1995 and 1997 solar minimum \cite{PhysRevLett.84.1078}), BESS Polar II (2007-08 solar minimum \cite{PhysRevLett.108.051102}), PAMELA (2006-09 with \~550 MV best-fit solar modulation potential \cite{Adriani2013}), and AMS-02 (2011-18 with average solar modulation potential ~620 MV \cite{PhysRevLett.117.091103, AGUILAR20211}) are also shown.
    \caption{The \ac{GAPS} antiproton sensitivity plot}
    \label{fig:GAPS_antiproton}
\end{figure}

\begin{figure}
    \centering
    \includegraphics[width=1\textwidth]{fig/GAPS_antideuteron_sensitivity.png}
    \noindent\raggedright The \ac{GAPS} antideuteron sensitivity (3$\sigma$ discovery for 1 detected event) for three LDB flights \cite{Xiao:2021V4}, in comparison with viable DM models \cite{PhysRevD.97.103011, Cui2010, BRAUNINGER200920, Hryczuk_2014, Randall2020}.
    \caption{The \ac{GAPS} antideuteron sensitivity plot}
    \label{fig:GAPS_antideuteron}
\end{figure}

We also estimated the \ac{aHe3} sensitivity of \ac{GAPS} which is also showed in Fig.~\ref{fig:cosmic_flux}. With or without an candidate events of \ac{aHe3} detected, the \ac{GAPS} data will provide a unique constraint on the cosmic \ac{aHe3} flux in the low energy range, which is crucial for understanding the origin of cosmic rays and their interactions with the interstellar medium.

As for \ac{GRAMS}, we have launched \ac{eGRAMS} as the world first ever \ac{LArTPC} detector in stratosphere. We are expecting to have a \ac{pGRAMS} launch in the coming month to verify the functionality of the instrument to have the first particle-tracking capability and compton camera \ac{LArTPC} detector ever launched on a balloon. The prototype flight \ac{pGRAMS}will also provide valuable data for optimizing the design of the full-scale \ac{GRAMS} instrument, which is expected to be launched in the next few years. With its unique capabilities, \ac{GRAMS} has the potential to make significant contributions to our understanding of $ MeV $ gamma-ray, cosmic rays and dark matter, see Fig.~\ref{fig:GRAMS_MeV_science}, Fig.~\ref{fig:cosmic_flux} and Fig.~\ref{fig:GRAMS_antiproton_excess}, and we are excited about the prospects of this experiment.

\begin{figure}
    \centering
    \includegraphics[width=1\textwidth]{fig/GRAMS_antiproton_exces.png}
    \noindent\raggedright Figure shows the parameter space for the possible dark matter signatures suggested by Fermi and AMS-02 \cite{PhysRevD.91.063003, DAYLAN20161, PhysRevD.93.083514, PhysRevLett.118.191101, 10.3847/1538-4357/aa6cab}, where GRAMS and GAPS sensitivities are overlaid with an uncertainty of the antideuteron production model \cite{PhysRevD.97.103011, N.Fornengo_2013, Acharya2017ProductionOD}.
    \caption{The galactic center excess\cite{aramaki2020dual}}
    \label{fig:GRAMS_antiproton_excess}
\end{figure}